\chapter{Introduction}
\label{chap:introduction}

Large extracellular electrophysiological recordings provide invaluable information about neuronal activity and interactions. However, for this information to be uncovered, a good spike sorting is necessary, assigning detected electrical events to putative neurons. In this process, numerous issues could potentially arise. The sorter could confuse neurons, miss spikes, or create spurious units. Such mistakes would naturally impact downstream scientific interpretation, for example, the firing rate and population structure. Classic validation work was done manually. This was not only tedious but also demonstrated that spike separation errors can vary significantly even under careful analysis. For over a decade now, the field has relied on quantitative quality metrics \parencite{doi:10.1152/jn.2000.84.1.401,Hill2011} and often a hard threshold to determine whether a sorted unit should be considered well-isolated or not.

The challenge lies in the fact that the most meaningful quality targets require ground-truth data for each neuron. Ground-truth data provide the information to identify \fpos{}, which measures the fraction of spikes assigned to a unit that should not be there (false positives), and \fmiss{}, which measures the fraction of true spikes the sorted unit did not capture (false negatives). In this thesis, both \fpos{} and \fmiss{} are targets we aim to predict. A model capable of doing this would be practically useful because it would assist in quantifying the risk of over-counting and under-counting neuronal activity. However, training such a model would require extensive ground-truth data, ideally using paired juxtacellular and intracellular ground-truth recordings that provide the most realistic ground-truth targets. Yet paired ground-truth recordings are rare and costly, and each recording only supplies a single ground-truth unit. Conversely, hybrid simulations are plentiful and comprehensively labelled, but they are not identical to real deployment data.

This asymmetry turns our prediction of \fpos{} and \fmiss{} into a transfer-learning challenge. The key question then becomes: can supervision from hybrid simulation data be effectively transferred to real paired ground-truth recordings? The main task is for the model to learn using hybrid simulations as our large labelled hybrid simulation dataset, then adapt to realistic noise and feature distributions using a small labelled paired ground-truth set plus a large unlabeled paired-recording pool, and finally evaluate under strict recording-disjoint protocols that mimic real operational use.

This thesis lies at the intersection of computational neuroscience, benchmark design, and low-label machine learning. Neuroscience allows us to identify the problem, define our features, and define the targets \fmiss{} and \fpos{}. Benchmarking practice determines what constitutes an honest test. Machine learning then allows us to compare different ways of leveraging the available datasets.

\chapter{Background}
\label{chap:background}

\section{History of Spike Sorting}

Monitoring the activity of specific neurons has been an area of research for over half a century. The first neuron activity was recorded using glass micropipettes and metal electrodes \parencite{Hubel1962}. At the time, identifying single spikes was a manual job performed by expert operators \parencite{Hubel1962}. Single-neuron activity can be detected with a single electrode and a simple threshold detector built into the hardware. However, the high level of background noise from other neurons in the local area, with action potentials of similar shape and size, made it challenging to isolate a specific neuron's activity \parencite{Lewicki1998}. The introduction of principal component analysis established it as the standard method for feature extraction in spike sorting. It marked the transition from manual methods to a software pipeline composed of four steps: detection, feature extraction, clustering, and evaluation \parencite{Lewicki1998}.

Since the beginning of extracellular recordings, our capacity to monitor multiple neurons simultaneously has doubled every seven years \parencite{Stevenson2011}. To capture more and more neural activity, the hardware had to evolve, starting from micropipettes and a singular metal electrode \parencite{Hubel1962} to modern high-density probes \parencite{Jun2017}. For the new generation of probes to achieve its full potential, it must be accompanied by the development of adapted spike-sorting algorithms and quality metrics \parencite{EINEVOLL201211}. With more neuronal activity recorded, additional challenges emerged: the signal became harder to parse, and classic sorting algorithms had to evolve to leverage the quantity of information provided by the new hardware.

\section{Post Hoc Analysis in Spike Sorting}

\subsection{Manual Sorting}

Spike sorting has traditionally depended on manual intervention and often still does. Experts inspect waveform clusters one channel at a time, in a two-dimensional PCA projection, and manually define cluster boundaries \parencite{Wood2004}. In an experiment led by \textcite{Wood2004}, five sorting experts were asked to perform manual spike sorting on 20 real motor-cortex channels and five realistic synthetic channels that experts could barely differentiate. The sorting experts agreed only 25\% of the time on the number of units in the real channel and had an average error rate of 23\% for false positives and 30\% for false negatives in synthetic channels \parencite{Wood2004}.

This experience demonstrated that manual spike sorting is highly variable, even amongst expert researchers, making it more difficult to draw conclusions from such data.

\subsection{Manual Curation}

As spike sorting became more automated, manual sorting disappeared. However, experts still perform manual curation tasks to improve performance, such as correcting the putative units identified by the sorters by merging split clusters, verifying cluster assignments, or rejecting poor-quality units \parencite{Rossant2016}.

However, manual curation is viewed as a limitation. It is not a scalable solution and it creates major reproducibility issues in the field \parencite{Rossant2016,Buccino2020SpikeInterface}. Recently, Fehér et al. (2025) proposed to automate unit curation using \emph{AECuration}, an autoencoder-based approach trained to separate true neural events from non-neural detections. Furthermore, Jain et al. (2025) introduced \emph{UnitRefine}, a community toolbox that uses supervised learning to classify putative units into three categories: noise, single-unit activity, and multi-unit activity. They trained the model using labels assigned by expert curators, whose predictions can vary and may deviate from the ground truth.

\subsection{Quality Metrics}

Robust quality metrics are vital for determining whether a unit identified by a spike-sorting algorithm corresponds to a single neuron. They are used by manual curators and scientists to assess the reliability of putative units. Quality metrics in spike sorting can be stratified into five categories: isolation, completeness, contamination, stability, and additional metrics.

Isolation metrics for clustering, such as isolation distance \parencite{Harris2001}, L-ratio \parencite{SchmitzerTorbert2005}, D-prime \parencite{Hill2011}, and silhouette score \parencite{ROUSSEEUW198753}, are used to quantify how well a putative unit is separated from other units and noise.

Contamination metrics, such as interspike interval violations \parencite{Hill2011}, noise cutoff, and signal-to-noise ratio \parencite{IBL2022}, measure whether the algorithm includes noise or spikes from other neurons in a cluster; in other words, they measure false positives.

Completeness metrics, such as amplitude cutoff \parencite{Hill2011}, amplitude median, amplitude coefficient of variation, and presence ratio \parencite{IBL2022}, measure whether the algorithm fails to see some spikes that actually belong to a neuron; in other words, they measure false negatives.

Drift metrics, such as drift peak-to-peak, drift standard deviation, and drift median absolute deviation \parencite{Siegle2021}, quantify the physical stability of a neuron relative to the recording electrodes.

Quality metrics are insightful but not sufficient to distinguish with perfect precision which putative units are accurate. In fact, they are proxies while the goal is to estimate the ground-truth error rate itself.

\section{Different Spike-Sorting Algorithms}

Spike-sorter families differ in their structure and in what they optimise for. Template-based methods are among the most popular: they use spike waveforms, which serve as the signature of a neuron, to determine the unit to which a spike belongs. From this group, we evaluate KiloSort \parencite{Pachitariu2016}, KiloSort2, SpyKING CIRCUS \parencite{10.7554/eLife.34518}, and Klusta \parencite{Rossant2016}. This approach particularly excels at resolving overlapping spikes and scales effectively to high-dimensional probes.

Similarly, JRClust \parencite{Jun2017} and its extension, IronClust, were designed for high-density probes, and their algorithms account for potential probe drift to prevent assignment errors. HerdingSpikes2, MountainSort4, and Tridesclous represent other important design choices, differing in their handling of dense arrays, automation, and practical integration. These spike sorters differ in terms of the probes they are designed for, the edge cases they consider, their robustness, their need for manual curation, and how they handle overlapping spikes. All of these factors influence how they sort spikes across different units. Therefore, a recording sorted by different sorters might yield different numbers of detected spikes or neurons.

\section{Machine Learning}

\subsection{Transfer Learning and Domain Shift}

Transfer learning is a machine-learning approach where knowledge gained from one task or domain is reused to enhance performance on a different, yet related, task or domain \parencite{PanYang2010}. When data are limited, transfer learning offers an excellent way to train a model on synthetic data and subsequently adapt it to real-world data. For instance, \textcite{Zhang2022-wg} describe their sim-to-real setting as related to supervised domain adaptation for few-shot learning, where only very few labelled target samples are available, and report strong performance. This motivates the idea that applying a similar transfer logic from synthetic spike-sorting data to real paired recordings could improve our model.

\subsection{Tabular Data Regression}

For small to medium tabular datasets, tree-based methods often outperform neural networks. \textcite{Grinsztajn2022} demonstrated that tree-based models remain state of the art on typical tabular datasets of this scale. XGBoost, introduced by \textcite{ChenGuestrin2016}, is particularly effective for small datasets because of its sparsity-aware algorithm, and LightGBM provides a similarly strong gradient-boosted-tree framework \parencite{Ke2017}.

\section{Research Gap}

Extensive research has been carried out to improve sorters and benchmark their performance. Meanwhile, we continue to assess the quality of putative units using proxy metrics. Some experiments have been conducted to identify noise units and multi-unit activity through human-curation data. However, classification does not fully capture small variations in unit quality, and even expert curators sometimes make mistakes, which can make the labels unreliable.

To the best of our knowledge, no previous work has attempted to predict the two continuous variables \fpos{} and \fmiss{}, which are part of the neuron-accuracy formula, using a supervised regression benchmark designed specifically around leakage-safe transfer from hybrid data to paired recordings.


\section{Research Questions and Hypotheses}

The thesis is organised around four hypotheses that reflect the most significant uncertainties, in \cref{tab:hypotheses}.

\begin{table}[H]
    \centering
    \caption{Research hypotheses introduced in the thesis.}
    \label{tab:hypotheses}
    \begin{tabularx}{\textwidth}{@{}lY@{}}
        \toprule
        Hypothesis & Statement \\
        \midrule
        H1 & Hybrid simulation data and paired juxtacellular and intracellular ground-truth data are strongly shifted, making naive hybrid-only transfer insufficient for reliable quality prediction. \\
        H2 & Leakage-safe recording context adds meaningful predictive signal. \\
        H3 & \fpos{} and \fmiss{} are structurally different prediction targets. \\
        H4 & The best model depends on the dataset family and evaluation regime rather than being universal across all paired recordings. \\
        \bottomrule
    \end{tabularx}
\end{table}

H1 questions the existence of a meaningful transfer problem that would require methods more complex than hybrid-only transfer. H2 tests whether information at the recording level can add significant signal. H3 and H4 move beyond performance and question the heterogeneity of our targets as well as whether different models perform better under different paired-recording families.

\section{Contributions}

The main contributions of the thesis are as follows.

\begin{itemize}
    \item It constructs a leakage-safe benchmark for predicting spike-sorting error metrics, using recording-disjoint and family-disjoint splits as its main evaluation method rather than mere unit-level random splits.
    \item It quantifies the shift between hybrid simulation data and paired juxtacellular and intracellular ground-truth data, showing near-perfect domain separability.
    \item It compares multiple families of transfer methods, including hybrid-only transfer, context-aware stacking, unlabeled paired-recording structure, and waveform-enriched stacking.
    \item It shows a clear target asymmetry: \fpos{} reveals a significant signal and meaningful transfer, while \fmiss{} shows greater instability.
    \item It demonstrates that family structure within paired ground-truth data matters. The best \fpos{} model differs across paired-recording families.
\end{itemize}

These contributions are deliberately scoped and do not claim to fully solve the quality-control problem.

\section{Why This Thesis Focuses on Quality Prediction Rather Than Spike Sorting End-to-End}

End-to-end spike sorting and post hoc quality estimation are related but distinct problems. First, from a noisy recording signal across multiple channels, neuronal activity, or spikes, needs to be identified. Then a spike sorter groups the detected spikes into sets of putative neurons, whose accuracy can range from perfectly accurate to fully erroneous, with spikes misassigned, missed, or attributed to spurious units. Between those two situations exists a continuous spectrum of intermediate levels of sorting accuracy. The quality predictor comes after that sorting stage to identify where each specific putative neuron lies on the spectrum.

State-of-the-art spike sorters have already drastically improved their performance over the last few decades, as they have been the main focus of the spike-sorting field. Post hoc quality estimation is not yet up to date with all the discoveries in both computational neuroscience and machine learning. This discrepancy highlights the necessity for more resilient, data-driven methods for post hoc quality evaluation. A reliable quality predictor would improve scientific analysis and curation and could later help us understand what constitutes a good neuron, which could support the next generation of spike-sorting algorithms.

\section{Roadmap}

The rest of the manuscript proceeds as follows. \Cref{chap:methodology} describes the data package, feature construction, model families, leakage controls, and evaluation protocols. \Cref{chap:results} presents the empirical results, including domain shift, model progression for \fpos{} and \fmiss{}, cross-target asymmetry, explainability, and label-budget constraints. \Cref{chap:discussion} returns to the hypotheses and interprets what the benchmark actually demonstrates. \Cref{chap:conclusion} closes with the main thesis-level takeaway.
